\section{Troubleshooting and practical recipes}
\label{sec:troubleshooting}

\subsection{Installation and environment}
\begin{description}[style=nextline]
\item[\texttt{ModuleNotFoundError: No module named 'pulsus'}]
Activate the intended environment and install with
\texttt{python -m pip install pulsus-spec}. Print
\texttt{sys.executable} and \texttt{pulsus.\_\_file\_\_} to identify the
interpreter and source actually used. Avoid naming your own script
\texttt{pulsus.py}, which can shadow the package.
\item[Notebook imports a different version.]
Select the kernel belonging to the spectroscopy environment. If needed,
install \texttt{ipykernel} in that environment and register it:
\par\noindent
\begin{minipage}{\linewidth}
\begin{lstlisting}[style=shell]
python -m pip install ipykernel
python -m ipykernel install --user --name pulsus --display-name "Python (PULSUS)"
\end{lstlisting}
\end{minipage}\par
Restart the selected kernel after updating a package already imported there.
\item[An example cannot import \texttt{dimer\_model} or \texttt{common}.]
The older development examples may use \texttt{src/dimer\_model.py}, which is
not an installed public PULSUS module. This manual's first example is
self-contained. Run examples~02--05 as scripts from the checkout, keeping
their \texttt{common.py} helper beside them; do not copy those scripts alone
and expect the helper to come from PyPI.
\end{description}

\subsection{States and operators}
\begin{description}[style=nextline]
\item[\texttt{An initial state must be supplied ...}]
Computing \texttt{rho\_ss = system.stationary\_state()} does not populate
\texttt{system.rho0}. Pass \texttt{rho0=rho\_ss} to the calculation, or assign
a suitable initial state when constructing the system. A call-level
\texttt{rho0} overrides the stored state.
\item[\texttt{RWA calculations require dipole\_plus and dipole\_minus ...}]
Supply both matrices to \texttt{SpectroscopySystem}. Check
$\mu=\mu_++\mu_-$ and, for a Hermitian dipole, $\mu_-=\mu_+^\dagger$.
The package does not infer the decomposition from $H$ automatically.
\item[Stationary state is ambiguous or appears unphysical.]
The eigenvector-based routine chooses the eigenvalue closest to zero. A
generator with multiple stationary states does not specify a unique
preparation. Check trace, Hermiticity, eigenvalues and residual, then supply
the state that represents the experiment. Adding arbitrary relaxation solely
to force uniqueness changes the physical model.
\item[Relaxation is much faster or slower than expected.]
Collapse operators contain the square root of a rate. Using $\gamma A$
where $\sqrt{\gamma}A$ was intended changes the dissipative rate to a quantity
proportional to $\gamma^2$. Also check that $H/\hbar$, rates, angular
frequencies and time parameters use compatible units.
\end{description}

\Needspace{12\baselineskip}
\subsection{Coordinates, amplitudes and plotting}
\begin{description}[style=nextline]
\item[The rephasing signal is absent on a positive excitation axis.]
The calculation uses negative \texttt{omega1} for the standard R sector.
\texttt{pathway="R"} selects interaction signs but does not flip coordinates.
Calculate on the negative axis, then display against $-\omega_1$ if desired.
\item[A spectrum appears transposed.]
Grid outputs have shape $(N_3,N_1)$, with detection on rows. Use unequal
axis lengths to diagnose orientation. Pair coordinates with
\texttt{finite\_pulse\_points}; do not reshape its $(N,)$ output into a
Cartesian grid unless you explicitly constructed and ordered that grid.
\item[Changing \texttt{pathway} has no effect.]
Check whether an explicit \texttt{signature} is also supplied. That tuple
overrides the pathway alias. Also distinguish full impulsive functions,
which do not have a pathway argument, from sector-selecting RWA functions.
\item[The plotted sign or amplitude is surprising.]
The response is complex. Choose \texttt{real}, \texttt{imag} or \texttt{abs}
deliberately and record the phase and normalization conventions. The plotter
defaults to the real component without normalization. Negative signed
components are valid. Normalizing each plot separately removes absolute
amplitude information and is unsuitable for amplitude-error comparisons.
\item[Very narrow pulses give a small signal.]
At fixed peak field, each Gaussian signed-component area scales as
$E_0\sigma$. Narrowing the pulse therefore also reduces its area. If the
purpose is an impulsive-limit comparison at fixed area, adjust $E_0$
consistently rather than comparing unlike field strengths.
\item[Values are nonfinite or unexpectedly large.]
Check input units, pulse widths, physical rates and stationarity first.
The Gaussian dressing contains matrix exponentials; extreme parameter
choices can expose numerical conditioning or fall outside the separated-pulse
regime. Start from the validated small example and change one parameter at a
time. Do not hide nonfinite values by normalizing or clipping the plot.
\end{description}

\subsection{Prepare a custom model}
Begin with a small truncated Hamiltonian, a dipole in the same basis and
physically justified collapse operators. Verify state preparation and a
linear spectrum before building a large third-order grid. For an electronic
or vibronic RWA calculation, define the raising/lowering decomposition in
the intended excitation manifolds. Then compare approximations on identical
coordinates using a documented pulse-amplitude convention.
